\documentclass[11pt]{article}

\usepackage[utf8]{inputenc}
\usepackage[T1]{fontenc}
\usepackage{amsmath, amssymb}
\usepackage{geometry}
\usepackage{setspace}
\usepackage{booktabs}
\usepackage{enumitem}
\usepackage{hyperref}

\geometry{a4paper, margin=2.5cm}
\onehalfspacing

\title{Plan de Extension del Framework PyExner para el Modelo EPB de Dos Fluidos}
\author{Esteban Vega Romero}
%\date{29 de abril de 2026}

\begin{document}

\maketitle

\section*{Resumen}

Este documento describe las fases de extension del framework \texttt{PyExner} para incorporar un modelo 2.5D de \emph{Burbujas de Plasma Ecuatoriales} (EPB) basado en una formulacion hiperb\'olica de dos fluidos. El objetivo no es reemplazar los solucionadores hidraulicos ya existentes, sino a\~nadir una nueva rama fisica \texttt{EPB\_TwoFluid} que conviva con los esquemas \texttt{Roe} y \texttt{Roe Exner}.

La idea central es reutilizar la infraestructura de ejecucion, malla, MPI, contornos, integracion temporal e I/O paralelo del framework, mientras se introduce un nuevo conjunto de variables conservadas, nuevos flujos hiperb\'olicos, un operador de fuentes separado y una etapa eliptica para el potencial de polarizacion.

El sistema continuo de interes se expresa como

\begin{equation}
\frac{\partial \mathbf{Q}}{\partial t}
+ \frac{\partial \mathbf{F}(\mathbf{Q})}{\partial x}
+ \frac{\partial \mathbf{G}(\mathbf{Q})}{\partial z}
= \mathbf{S}(\mathbf{Q}),
\end{equation}

donde el vector de estado conservado es

\begin{equation}
\mathbf{Q} =
\begin{bmatrix}
n_i & n_e & j_{ix} & j_{iy} & j_{iz} & j_{ex} & j_{ey} & j_{ez}
\end{bmatrix}^{T}.
\end{equation}

Una restriccion de dise\~no fundamental es mantener separadas las siguientes piezas numericas:

\begin{enumerate}[label=\arabic*.]
    \item el transporte hiperb\'olico explicitamente conservativo,
    \item el operador de fuentes rigidas de colision y electrodinamica,
    \item el solve eliptico que recupera el campo electrico de polarizacion mediante $\mathbf{E} = \mathbf{E}_0 - \nabla \phi$.
\end{enumerate}

\section{Principios de Diseño}

La extension del framework se guiara por los siguientes principios:

\begin{enumerate}[label=\arabic*.]
    \item \textbf{Extension aditiva del framework.} Los modelos hidraulicos existentes deben seguir funcionando sin cambios de comportamiento.
    \item \textbf{Separacion entre fisica e infraestructura.} La arquitectura de \texttt{PyExner} se conserva; lo que cambia es la fisica implementada en una nueva rama de solver.
    \item \textbf{Separacion entre flujo, fuentes y restricciones.} El paso hiperb\'olico no debe mezclar la dinamica de transporte con colisiones ni con el solve de $\phi$.
    \item \textbf{Consistencia del estado conservado.} El orden de variables en memoria, registros, kernels y pruebas debe coincidir exactamente con el orden fisico de $\mathbf{Q}$.
    \item \textbf{Implementacion vectorizada.} Los kernels numericos deben expresarse con \texttt{jax.numpy}, sin bucles Python por celda.
\end{enumerate}

\section{Fase 1: Generalizacion Minima del Framework}

\subsection*{Objetivo}

Eliminar las suposiciones estrictamente hidraulicas del framework compartido, pero sin modificar la API externa de los modelos existentes.

\subsection*{Motivacion}

En el estado actual, algunas piezas comunes del framework estan escritas bajo la hipotesis de que el estado siempre contiene variables como $h$, $hu$ y $hv$. Para EPB, el estado conservado tiene ocho componentes distintas y no debe forzarse dentro de nombres hidraulicos artificiales.

\subsection*{Trabajo tecnico}

\begin{enumerate}[label=\arabic*.]
    \item Generalizar las utilidades de estado para que operen sobre cualquier \texttt{dataclass} de arreglos.
    \item Revisar los \emph{NamedTuple} y tipos del integrador temporal para aceptar un estado arbitrario compatible con PyTrees de JAX.
    \item Mantener compatibilidad hacia atras con \texttt{Roe} y \texttt{Roe Exner}.
\end{enumerate}

\subsection*{Entregable}

Una base de framework capaz de transportar estados no hidraulicos sin romper la funcionalidad ya existente.

\subsection*{Riesgo principal}

El riesgo en esta fase es introducir una generalizacion demasiado agresiva y afectar los modelos SWE ya verificados. Por eso los cambios deben ser minimos y localizados.

\section{Fase 2: Alta de la Rama EPB\_TwoFluid}

\subsection*{Objetivo}

Crear la nueva rama fisica \texttt{EPB\_TwoFluid} dentro del sistema de registros del framework.

\subsection*{Trabajo tecnico}

\begin{enumerate}[label=\arabic*.]
    \item Crear un nuevo estado conservado dedicado para EPB.
    \item Registrar el nuevo estado con la cadena exacta \texttt{EPB\_TwoFluid}.
    \item Crear y registrar el nuevo \emph{solver bundle} con la misma cadena.
    \item Activar los imports necesarios para que el registro ocurra automaticamente al cargar el paquete.
\end{enumerate}

\subsection*{Razon numerica}

La formulacion EPB no es una variacion menor del solver de Roe. El modelo tiene nuevas variables, nuevos flujos y nuevas escalas temporales. Por tanto, debe implementarse como una rama fisica propia y no como un caso especial dentro del codigo hidraulico.

\subsection*{Entregable}

Un nuevo camino funcional dentro del framework:

\begin{equation}
\texttt{flux\_scheme: EPB\_TwoFluid}
\Longrightarrow
\texttt{state} + \texttt{solver} + \texttt{boundary} + \texttt{integrator}.
\end{equation}

\section{Fase 3: Slice Hiperbolico Puro}

\subsection*{Objetivo}

Implementar el nucleo de transporte hiperb\'olico del sistema EPB de dos fluidos sin incluir todavia las fuentes rigidas ni el solve electrostatico.

\subsection*{Contenido matematico}

Se implementaran los tensores de flujo $\mathbf{F}(\mathbf{Q})$ y $\mathbf{G}(\mathbf{Q})$ con la convencion de signos correcta para electrones e iones. Esto es especialmente delicado en los terminos de flujo electronico. Por ejemplo, el flujo en $x$ para el momento electronico en la direccion $x$ debe conservar la estructura

\begin{equation}
F_6 = -\frac{j_{ex}^{2}}{e n_e} - \frac{e}{M_e} p_e,
\end{equation}

y su permutacion correspondiente en la direccion vertical debe preservarse en $\mathbf{G}(\mathbf{Q})$.

\subsection*{Trabajo tecnico}

\begin{enumerate}[label=\arabic*.]
    \item Implementar kernels de flujo con \texttt{jax.numpy}.
    \item Calcular cotas de velocidad de onda para un esquema tipo HLL o HLLC.
    \item Construir un \emph{step} hiperb\'olico sin reusar eigensistemas de SWE.
    \item Calcular $\Delta t$ con base en la velocidad maxima de onda y reduccion MPI global.
\end{enumerate}

\subsection*{Entregable}

Un solver de transporte que preserve estados uniformes y que pueda ejecutarse con condiciones iniciales simples sin activar aun los terminos fuente.

\subsection*{Riesgo principal}

El mayor riesgo es reutilizar por comodidad formulas de Roe/SWE y contaminar el modelo EPB con autovalores o estructuras de onda que no le pertenecen.

\section{Fase 4: Contornos e I/O para EPB}

\subsection*{Objetivo}

Integrar el nuevo modelo EPB con las piezas de entrada, salida y manejo de frontera del framework.

\subsection*{Trabajo tecnico}

\begin{enumerate}[label=\arabic*.]
    \item Definir condiciones de contorno transmisivas para los ocho campos del estado EPB.
    \item Asegurar que la convencion de registro use nombres del tipo
    \begin{equation}
    \texttt{EPB\_TwoFluid Transmissive}.
    \end{equation}
    \item Extender la lectura NetCDF para poblar las ocho variables conservadas.
    \item Extender la escritura NetCDF para exportar esas variables sin asumir estructura hidraulica.
\end{enumerate}

\subsection*{Observacion importante}

La capa de I/O de \texttt{PyExner} ya tiene bastante generalidad porque recorre dinamicamente los campos del estado. Por eso, el trabajo principal no sera rehacer la clase de I/O, sino fijar de forma consistente los nombres, shapes y convenciones de entrada/salida para EPB.

\subsection*{Entregable}

Un caso EPB capaz de leer condiciones iniciales, aplicar contornos y escribir resultados sin modificaciones manuales en el pipeline de ejecucion.

\section{Fase 5: Fuentes Rigidas y Acoplamiento Electrost\'atico}

\subsection*{Objetivo}

Incorporar el operador de fuentes $\mathbf{S}(\mathbf{Q})$ y la etapa eliptica asociada al potencial de polarizacion $\phi$.

\subsection*{Contenido fisico}

El vector fuente reune al menos los siguientes mecanismos:

\begin{enumerate}[label=\arabic*.]
    \item gravedad,
    \item fuerza electrica,
    \item fuerza magnetica,
    \item colisiones ion-neutro,
    \item colisiones electron-neutro,
    \item colisiones electron-ion.
\end{enumerate}

El campo electrico efectivo se actualiza segun

\begin{equation}
\mathbf{E} = \mathbf{E}_0 - \nabla \phi,
\end{equation}

y la restriccion de conservacion de carga se recupera a traves de la ecuacion eliptica asociada al potencial, en lugar de imponerse dentro del paso hiperb\'olico.

\subsection*{Trabajo tecnico}

\begin{enumerate}[label=\arabic*.]
    \item Implementar el operador separado de fuentes en un kernel propio.
    \item Definir con claridad el orden de actualizacion entre transporte, fuentes y solve de $\phi$.
    \item Preparar la estructura para parametros de fondo: $g$, $E_x$, $E_y$, $E_z$, $B_0$, $U_x$, $U_y$, $U_z$, $\nu_{in}$, $\nu_{en}$ y $\nu_{ei}$.
\end{enumerate}

\subsection*{Riesgo principal}

Si las fuentes rigidas se integran de manera completamente explicita dentro del mismo paso del transporte, el esquema puede volverse numericamente inestable aun cuando el CFL convectivo sea pequeno.

\section{Fase 6: Integrador IMEX}

\subsection*{Objetivo}

Introducir un integrador de tipo IMEX que trate de forma explicita el transporte y de forma implicita las fuentes locales rigidas.

\subsection*{Razon numerica}

El paso hiperb\'olico esta gobernado por velocidades de onda, mientras que la parte fuente puede estar controlada por frecuencias colisionales mucho mas rapidas. Un unico integrador explicito obligaria a un paso temporal excesivamente pequeno o directamente inestable.

\subsection*{Trabajo tecnico}

\begin{enumerate}[label=\arabic*.]
    \item Crear un nuevo integrador \texttt{IMEX} como rama adicional del framework.
    \item Reutilizar el mismo contrato de ejecucion del framework: configuracion, bucle temporal y escritura de salida.
    \item Resolver la parte implicita con algebra local por celda, sin esconder la rigidez dentro del flujo hiperb\'olico.
\end{enumerate}

\subsection*{Entregable}

Un integrador estable para EPB que permita acoplar transporte, fuentes y solve electrostatico sin destruir la estructura del framework.

\section{Fase 7: Validacion}

\subsection*{Objetivo}

Verificar, en cada etapa, que la extension EPB no contiene errores estructurales de registro, signos, orden de variables o acoplamiento MPI.

\subsection*{Bateria minima de validacion}

\begin{enumerate}[label=\arabic*.]
    \item \textbf{Wiring del framework:} comprobar que \texttt{EPB\_TwoFluid} crea el estado y el solver correctos.
    \item \textbf{Orden del estado:} verificar que el empaquetado y desempaquetado de $\mathbf{Q}$ respeta el orden fisico definido.
    \item \textbf{Signo del flujo electronico:} validar especificamente el signo de la presion y de la continuidad electronica.
    \item \textbf{Contornos:} confirmar que el registro de fronteras sigue la convencion del framework.
    \item \textbf{Estado uniforme:} un estado constante debe permanecer constante en ausencia de fuentes y gradientes.
    \item \textbf{Halo exchange:} verificar que las ocho variables requeridas se sincronizan entre subdominios.
\end{enumerate}

\subsection*{Entregable}

Un conjunto pequeno pero discriminante de pruebas que detecte errores catastroficos temprano, antes de abordar simulaciones fisicas complejas.

\section{Orden de Ejecucion Recomendado}

El orden recomendado de implementacion es el siguiente:

\begin{enumerate}[label=\arabic*.]
    \item Generalizacion minima del framework compartido.
    \item Alta del estado \texttt{EPB\_TwoFluid} y su solver bundle.
    \item Implementacion del slice hiperb\'olico puro.
    \item Integracion de contornos e I/O.
    \item Implementacion del operador de fuentes y del solve de $\phi$.
    \item Incorporacion del integrador IMEX.
    \item Validacion incremental y casos de prueba canonicos.
\end{enumerate}

Este orden minimiza el riesgo de mezclar problemas de infraestructura con problemas de fisica o estabilidad numerica.

\section{Conclusiones}

La extension EPB de \texttt{PyExner} es viable siempre que se conserve una separacion estricta entre infraestructura compartida y fisica especifica del nuevo modelo. La reutilizacion debe concentrarse en la arquitectura del framework, no en el nucleo numerico de los solucionadores hidraulicos.

Desde el punto de vista de ingenieria de software, el proyecto consiste en convertir \texttt{PyExner} en un framework con una nueva familia de solucionadores, sin alterar la validez de las ramas \texttt{Roe} y \texttt{Roe Exner}. Desde el punto de vista numerico, el requisito critico es mantener desacoplados el transporte hiperb\'olico, los terminos fuente rigidos y el solver electrostatico.

\end{document}